% Generated draft: requires mathematical and source review.
\documentclass[11pt]{article}
\usepackage[margin=25mm]{geometry}
\usepackage{fontspec}
\setmainfont{Latin Modern Roman}
\usepackage{amsmath,amssymb,mathtools,unicode-math}
\setmathfont{Latin Modern Math}
\usepackage{enumitem,needspace,xurl,hyperref}
\setlength{\parindent}{0pt}
\setlength{\parskip}{5pt}
\setlength{\emergencystretch}{3em}
\newcommand{\N}{\mathbb{N}}
\newcommand{\Z}{\mathbb{Z}}
\newcommand{\Q}{\mathbb{Q}}
\newcommand{\R}{\mathbb{R}}
\newcommand{\C}{\mathbb{C}}
\newcommand{\F}{\mathbb{F}}
\newcommand{\A}{\mathbb{A}}
\newcommand{\PP}{\mathbb{P}}
\newcommand{\E}{\mathbb{E}}
\newcommand{\eps}{\varepsilon}
\newcommand{\dd}{\,\mathrm{d}}
\begin{document}
\section*{UnsolvedMath ID: 800077}
\textbf{Source problem code:} catalogue entry
\subsection{Definitions and mathematical statement}
For a smooth bundle \(\pi\colon E\to B\) with closed, connected, oriented surface fibre \(\Sigma_g\) and orientation-preserving transition maps, let \(T_\pi E\) be its oriented vertical tangent bundle. Its Miller--Morita--Mumford classes are
\[
\kappa_i(\pi)=\pi_!\bigl(e(T_\pi E)^{i+1}\bigr)\in H^{2i}(B;\R),
\]
where \(e\) is the real Euler class and \(\pi_!\) is integration over the oriented surface fibre.

The source writes \(G=\mathrm{Diff}(\Sigma_g)\). It does not specify the genus range, whether diffeomorphisms preserve orientation, or whether this group is given the discrete topology. Under the interpretation that \(G\) is the discrete group of orientation-preserving diffeomorphisms, the preceding construction on the universal associated surface bundle gives \(\kappa_i\in H^{2i}(G;\R)\).

For a discrete group \(G\), let \(H_b^n(G;\R)\) denote the cohomology of bounded group cochains with trivial real coefficients, and let
\[
c_G\colon H_b^n(G;\R)\longrightarrow H^n(G;\R)
\]
be the comparison map induced by inclusion into all group cochains. The source attributes to Morita, using Bott vanishing, the assertion that \(\kappa_i=0\) in ordinary real cohomology for \(i>2\).

\textbf{Question, conditional on specifying the intended bounded classes.} If \(\kappa_i^b\in H_b^{2i}(G;\R)\) denotes the bounded MMM class intended by the source, with \(c_G(\kappa_i^b)=\kappa_i\), does
\[
\kappa_i^b=0\quad\text{in }H_b^{2i}(G;\R)
\qquad\text{for every integer }i>2
\]
hold for the intended genera?

The supplied record does not construct or identify \(\kappa_i^b\). The comparison identity alone does not specify such a class: once \(\kappa_i=0\), the zero bounded class is already a lift. Thus merely asking whether a vanishing bounded lift exists would not capture the intended question. The intended bounded classes, genus quantifiers, and group conventions require clarification from the source.
\subsection{Short English statement}
The source asks whether the higher MMM classes, whose ordinary real-cohomology vanishing it attributes to Morita, also vanish as bounded-cohomology classes. A precise formulation requires identification of the intended bounded MMM classes and specification of the genus and diffeomorphism-group conventions.
\subsection{Sources}
\begin{itemize}
\item Local UnsolvedMath registry: permanent ID \(800077\), exact source code \texttt{catalogue entry}, title ``Foliations and Diffeomorphism Groups --- Question 18.'' The supplied registry labels the problem open; this status is not independently verified.
\item Nariman and Yazdi, \emph{Problem List on Foliations and Diffeomorphism Groups} (2025), Question 18. Supplied source URL: \url{https://arxiv.org/abs/2503.18274}. The provenance records access on 2026-10-16 and extraction method \texttt{arxiv\_latex}; these are supplied provenance claims, not an assertion of independent access or verification.
\item The supplied statement attributes ordinary-cohomology vanishing for \(i>2\) to Morita using the Bott vanishing theorem. No separate bibliographic reference for Morita's result is supplied.
\end{itemize}
\end{document}
